\section{Variable-density results}
\label{sec:variable_density}

We now consider the full variable-density formulation. To isolate the effects of density and
viscosity stratification, we fix \(\Sch=1\) and vary \(\Atw\) and \(\Amu\). We first group the
weakly and moderately stratified cases, \(\Atw=0.05\) and \(0.5\), and compare the three viscosity
contrasts \(\Amu=-0.9,0,0.9\). The strongly stratified case \(\Atw=0.75\) is treated separately.

\subsection{Growth-rate maps across the parameter sweep}
\label{sec:vd_sigma_maps}

Before isolating individual effects, we survey the per-wavenumber mean growth rate
\(\overline{\sigma}_{1D}(k,t)\) over the sweep. \Cref{fig:VD_sigma_map_Sc0p25,fig:VD_sigma_map_Sc1,fig:VD_sigma_map_Sc10}
report it at \(\Sch=0.25\), \(1\), and \(10\), each covering three density contrasts and three
viscosity contrasts at fixed \(\delta_0=1\). The \(\Sch=0.25\) case evolves an order of magnitude
more slowly and amplifies roughly half as strongly as the others, so it is shown over a longer
window and with a tighter colour range; the remaining two share a common scale.

\begin{figure}[p]
  \centering
  \begin{minipage}{0.32\textwidth}
    \centering
    \tikzfig{sigma_map_VD_Sc_0p25_d0_1_At_0p1_Amu_m0p9}
  \end{minipage}
  \hfill
  \begin{minipage}{0.32\textwidth}
    \centering
    \tikzfig{sigma_map_VD_Sc_0p25_d0_1_At_0p1_Amu_0}
  \end{minipage}
  \hfill
  \begin{minipage}{0.32\textwidth}
    \centering
    \tikzfig{sigma_map_VD_Sc_0p25_d0_1_At_0p1_Amu_0p9}
  \end{minipage}
  \\[2pt]
  \begin{minipage}{0.32\textwidth}
    \centering
    \tikzfig{sigma_map_VD_Sc_0p25_d0_1_At_0p5_Amu_m0p9}
  \end{minipage}
  \hfill
  \begin{minipage}{0.32\textwidth}
    \centering
    \tikzfig{sigma_map_VD_Sc_0p25_d0_1_At_0p5_Amu_0}
  \end{minipage}
  \hfill
  \begin{minipage}{0.32\textwidth}
    \centering
    \tikzfig{sigma_map_VD_Sc_0p25_d0_1_At_0p5_Amu_0p9}
  \end{minipage}
  \\[2pt]
  \begin{minipage}{0.32\textwidth}
    \centering
    \tikzfig{sigma_map_VD_Sc_0p25_d0_1_At_0p75_Amu_m0p9}
  \end{minipage}
  \hfill
  \begin{minipage}{0.32\textwidth}
    \centering
    \tikzfig{sigma_map_VD_Sc_0p25_d0_1_At_0p75_Amu_0}
  \end{minipage}
  \hfill
  \begin{minipage}{0.32\textwidth}
    \centering
    \tikzfig{sigma_map_VD_Sc_0p25_d0_1_At_0p75_Amu_0p9}
  \end{minipage}
  \tikzfig{sigma_map_VD_colorbar_Sc_0p25}
  \caption{Per-wavenumber mean growth rate \(\overline{\sigma}_{1D}(k,t)\) in the
  variable-density regime at \(\Sch=0.25\) and \(\delta_0=1\). Rows correspond to
  \(\Atw=0.1\), \(0.5\), and \(0.75\); columns correspond to \(\Amu=-0.9\), \(0\),
  and \(0.9\), as labelled above each panel. The (\solidlegend) contour marks neutral
  growth, \(\overline{\sigma}_{1D}=0\), and the (\dashlegend) curve identifies the most
  amplified wavenumber \(k_\ast(t)\). Each panel insets the early-time window
  \(t\leq8\). Color limits are \(\overline{\sigma}_{1D}\in\pm0.2\), shared across every
  panel of this figure.}
  \label{fig:VD_sigma_map_Sc0p25}
\end{figure}

\begin{figure}[p]
  \centering
  \begin{minipage}{0.32\textwidth}
    \centering
    \tikzfig{sigma_map_VD_Sc_1_d0_1_At_0p1_Amu_m0p9}
  \end{minipage}
  \hfill
  \begin{minipage}{0.32\textwidth}
    \centering
    \tikzfig{sigma_map_VD_Sc_1_d0_1_At_0p1_Amu_0}
  \end{minipage}
  \hfill
  \begin{minipage}{0.32\textwidth}
    \centering
    \tikzfig{sigma_map_VD_Sc_1_d0_1_At_0p1_Amu_0p9}
  \end{minipage}
  \\[2pt]
  \begin{minipage}{0.32\textwidth}
    \centering
    \tikzfig{sigma_map_VD_Sc_1_d0_1_At_0p5_Amu_m0p9}
  \end{minipage}
  \hfill
  \begin{minipage}{0.32\textwidth}
    \centering
    \tikzfig{sigma_map_VD_Sc_1_d0_1_At_0p5_Amu_0}
  \end{minipage}
  \hfill
  \begin{minipage}{0.32\textwidth}
    \centering
    \tikzfig{sigma_map_VD_Sc_1_d0_1_At_0p5_Amu_0p9}
  \end{minipage}
  \\[2pt]
  \begin{minipage}{0.32\textwidth}
    \centering
    \tikzfig{sigma_map_VD_Sc_1_d0_1_At_0p75_Amu_m0p9}
  \end{minipage}
  \hfill
  \begin{minipage}{0.32\textwidth}
    \centering
    \tikzfig{sigma_map_VD_Sc_1_d0_1_At_0p75_Amu_0}
  \end{minipage}
  \hfill
  \begin{minipage}{0.32\textwidth}
    \centering
    \tikzfig{sigma_map_VD_Sc_1_d0_1_At_0p75_Amu_0p9}
  \end{minipage}
  \tikzfig{sigma_map_VD_colorbar_Sc_1}
  \caption{Per-wavenumber mean growth rate \(\overline{\sigma}_{1D}(k,t)\) in the
  variable-density regime at \(\Sch=1\) and \(\delta_0=1\). Rows correspond to
  \(\Atw=0.1\), \(0.5\), and \(0.75\); columns correspond to \(\Amu=-0.9\), \(0\),
  and \(0.9\), as labelled above each panel. The (\solidlegend) contour marks neutral
  growth, \(\overline{\sigma}_{1D}=0\), and the (\dashlegend) curve identifies the most
  amplified wavenumber \(k_\ast(t)\). Each panel insets the early-time window
  \(t\leq2\). Color limits are \(\overline{\sigma}_{1D}\in\pm0.5\), shared across every
  panel of this figure.}
  \label{fig:VD_sigma_map_Sc1}
\end{figure}

\begin{figure}[p]
  \centering
  \begin{minipage}{0.32\textwidth}
    \centering
    \tikzfig{sigma_map_VD_Sc_10_d0_1_At_0p1_Amu_m0p9}
  \end{minipage}
  \hfill
  \begin{minipage}{0.32\textwidth}
    \centering
    \tikzfig{sigma_map_VD_Sc_10_d0_1_At_0p1_Amu_0}
  \end{minipage}
  \hfill
  \begin{minipage}{0.32\textwidth}
    \centering
    \tikzfig{sigma_map_VD_Sc_10_d0_1_At_0p1_Amu_0p9}
  \end{minipage}
  \\[2pt]
  \begin{minipage}{0.32\textwidth}
    \centering
    \tikzfig{sigma_map_VD_Sc_10_d0_1_At_0p5_Amu_m0p9}
  \end{minipage}
  \hfill
  \begin{minipage}{0.32\textwidth}
    \centering
    \tikzfig{sigma_map_VD_Sc_10_d0_1_At_0p5_Amu_0}
  \end{minipage}
  \hfill
  \begin{minipage}{0.32\textwidth}
    \centering
    \tikzfig{sigma_map_VD_Sc_10_d0_1_At_0p5_Amu_0p9}
  \end{minipage}
  \\[2pt]
  \begin{minipage}{0.32\textwidth}
    \centering
    \tikzfig{sigma_map_VD_Sc_10_d0_1_At_0p75_Amu_m0p9}
  \end{minipage}
  \hfill
  \begin{minipage}{0.32\textwidth}
    \centering
    \tikzfig{sigma_map_VD_Sc_10_d0_1_At_0p75_Amu_0}
  \end{minipage}
  \hfill
  \begin{minipage}{0.32\textwidth}
    \centering
    \tikzfig{sigma_map_VD_Sc_10_d0_1_At_0p75_Amu_0p9}
  \end{minipage}
  \tikzfig{sigma_map_VD_colorbar_Sc_10}
  \caption{Per-wavenumber mean growth rate \(\overline{\sigma}_{1D}(k,t)\) in the
  variable-density regime at \(\Sch=10\) and \(\delta_0=1\). Rows correspond to
  \(\Atw=0.1\), \(0.5\), and \(0.75\); columns correspond to \(\Amu=-0.9\), \(0\),
  and \(0.9\), as labelled above each panel. The (\solidlegend) contour marks neutral
  growth, \(\overline{\sigma}_{1D}=0\), and the (\dashlegend) curve identifies the most
  amplified wavenumber \(k_\ast(t)\). Each panel insets the early-time window
  \(t\leq2\). Color limits are \(\overline{\sigma}_{1D}\in\pm0.5\), shared across every
  panel of this figure.}
  \label{fig:VD_sigma_map_Sc10}
\end{figure}

\subsection{Low-to-moderate stratification}
\label{sec:vd_low_moderate}

As in the Boussinesq analysis, we begin with the per-wavenumber mean growth rate
\(\overline{\sigma}_{1D}(k,t)\) defined in~\eqref{eq:sigma_1d}. \Cref{fig:VD_P1_sigma_map}
shows its evolution for the six weakly and moderately stratified cases.

\begin{figure}[p]
  \centering
  \tikzfig{VD_P1_Sc1_At005_At050_sigma_grid}
  \caption{Per-wavenumber mean growth rate \(\overline{\sigma}_{1D}(k,t)\) in the
  variable-density regime at \(\Sch=1\). The first and second rows correspond to
  \(\Atw=0.05\) and \(0.5\), respectively; the columns correspond to \(\Amu=-0.9\), \(0\), and
  \(0.9\). The (\solidlegend) contour marks neutral growth, \(\overline{\sigma}_{1D}=0\), and the (\dashlegend)
  curve identifies the most amplified wavenumber \(k_\ast(t)\). Color limits are shared within
  each row.}
  \label{fig:VD_P1_sigma_map}
\end{figure}

Across all six cases in \cref{fig:VD_P1_sigma_map}, the evolution follows the same broad
sequence found in the Boussinesq limit. The initial response is stable, the growth-rate field then
crosses neutrality and reaches a positive maximum, and the instability subsequently weakens as
the diffusive layer broadens. Finite density and viscosity contrasts therefore modify the strength
of the response and the selected horizontal scale without changing this qualitative sequence.

The effect of viscosity contrast is most apparent in the dashed locus marking \(k_\ast(t)\).
At fixed \(\Atw\), increasing \(\Amu\) shifts this locus towards larger wavenumbers and increases
the maximum value reached by \(k_\ast(t)\). Because \(\Atw>0\) identifies fluid 2 as the heavier
fluid, \(\Amu<0\) implies that the lighter fluid is more viscous. The stronger viscous damping on
the lighter-fluid side inhibits small-scale motion, leading to a smaller peak value of
\(k_\ast(t)\) and hence a larger dominant horizontal scale, proportional to \(k_\ast^{-1}\).
For \(\Amu>0\), the viscosity ordering is reversed. Smaller scales are less strongly damped, and
the selected wavenumber reaches a larger maximum. The upward displacement of the neutral contour
with increasing \(\Amu\) similarly shows that the unstable band extends to larger wavenumbers.

The magnitude of the growth rate follows the same ordering once the neutral contour has been
crossed. Relative to \(\Amu=0\), a negative viscosity contrast reduces the growth rate throughout
the unstable stage, whereas a positive contrast produces larger values before the subsequent
decay. This behaviour is consistent between the weakly stratified case \(\Atw=0.05\) and the
moderately stratified case \(\Atw=0.5\), indicating that viscosity contrast controls both scale
selection and amplification over this range of density contrasts.

The value of the growth-rate field along the dashed locus defines
\(\overline{\sigma}_\ast(t)=\overline{\sigma}_{1D}(k_\ast(t),t)\), consistently with
\eqref{eq:maximum_growth_definitions}. We now test the proposed Atwood-number scaling for the
selected wavenumber,
\begin{equation}
  k_\ast(t)\sim t^{-(1-\Atw)/2}.
  \label{eq:vd_kstar_scaling}
\end{equation}
\Cref{fig:VD_P2_scaling} shows the compensated quantity for the same six cases considered
above.

\begin{figure}[p]
  \centering
  \begin{minipage}{0.48\textwidth}
    \centering
    \tikzfig{VD_P2_Sc1_At005_scaling}
  \end{minipage}
  \hfill
  \begin{minipage}{0.48\textwidth}
    \centering
    \tikzfig{VD_P2_Sc1_At050_scaling}
  \end{minipage}
  \caption{Compensated most amplified wavenumber \(k_\ast(t)t^{(1-\Atw)/2}\) at \(\Sch=1\).
  Panels (a) and (b) correspond to \(\Atw=0.05\) and \(0.5\), respectively, and each compares
  \(\Amu=-0.9,0,0.9\). An approximately flat interval identifies the time
  range over which the scaling proposed in~\eqref{eq:vd_kstar_scaling} holds; no collapse is
  implied outside that interval. The corresponding unscaled \(k_\ast(t)\) curves are the dashed
  loci in \cref{fig:VD_P1_sigma_map}.}
  \label{fig:VD_P2_scaling}
\end{figure}
